\chapter{Prägungen}
\label{app:marks}


\section{Abstammungen}\label{sec:abstammungen}

{
    \begin{multicols*}{2}

        \subsection{Menschen}\label{subsec:abstammungen-menschen}

        Die menschlichen Reiche Tirakans kamen zu Beginn des vierten Zeitalters auf, und ihre Blüte sollten sie zum Ende dieses Zeitalters hin erlangen. Die Reiche der Menschen schmiegen sich um das Felsenmeer, welches das Zentrum des Kontinentes bildet. Kontakt zu den anderen Kulturen Tirakans sollten die Menschen erst im Laufe der ersten Jahrhunderte des aktuellen Zeitalters erlangen.

        \subsubsection*{Al Bah JiRa}
        Im Land der Sonne, Al Bah JiRa, bestimmen uralte Glaubensvorstellungen, Stammesfehden und der stetige Kampf zwischen Einheit und Zerfall seit Jahrtausenden das Schicksal eines stolzen Wüstenvolkes. Aus der Offenbarung des Karawanenführers Apademak entstand einst die JiRa, die Lehre von Sonne und Wasser, welche die zerstrittenen Stämme einen sollte. Doch Bürgerkriege, alte Herrscherdynastien, verbotene Kulte und äußere Bedrohungen führten das Reich immer wieder an den Rand des Untergangs. Zwischen heiligen Oasen, mächtigen Hafenstädten und düsteren Geheimnissen bewahren die Menschen Al Bah JiRas bis heute ihre alte Kultur, ihre Geschichten und ihren Glauben.

        \abstammungspielwertebox
        {Erhalte +10 auf Anrufungen des Glaubens und formale Ritualführung.}
        {Wenn du einen geschworenen Ritus oder ein Tabu brichst, kostet dich das sofort 1 Omen (oder +1 Bürde bei 0 Omen).}

        \subsubsection*{Asgoran}
        Das Volk des Königreiches Asgoran ist ein Volk von Händlern und Seefahrern. Direkt am Rande des östlichen Ozeans gelegen herrscht hier oft eine steife Brise, obgleich das Klima ansonsten eher gemäßigt ist. Das Volk der Asgoraner besteht seit etwa 500 v.EC und hat aufgrund seines vielfältigen Handels Kontakt zu fast allen bekannten Reichen. Die Sprache der Asgoraner, Asgal ist gleichermaßen auch die Sprache, welche nahezu überall auf Tirakan verstanden wird.

        \abstammungspielwertebox
        {Erhalte +10 auf Proben mit Geschick + Athletik oder Geist + Mythen, die mit Wasser, Schiffen oder Küsten zu tun haben.}
        {In engem Inlandsgelände (enge Gassen, tiefe Höhlen) erhältst du -10 auf Bewegungsproben bei Verfolgung oder Flucht.}

        \subsubsection*{Die Stämme der Barbaren}
        Die Barbaren-Stämme sind die südlichsten Vertreter der Menschen auf Tirakan. Gespalten in acht verschiedene Stämme ist ihr Antlitz und ihr Verhalten höchst unterschiedlich. Die Barbaren des Südens waren die ersten, auf die der Zug der Minotauren mit voller Wucht prallte. So ergab es sich, dass bereits einige der Stämme im dritten Jahrhundert völlig ausgelöscht sind. Andere Stämme leben nach wie vor in den Landen südlich des Blutwalls neben den Minotauren.

        \abstammungspielwertebox
        {Erhalte +10 auf Proben zu Überleben, Fährtensuche, Kampf unter widrigen Bedingungen und dem Einschätzen fremder Bedrohungen.}
        {Jahrelanger Kampf und Stammesdenken nähren tiefes Misstrauen. Erhalte -10 auf die erste soziale Probe gegenüber fremden Kulturen oder unbekannten Autoritäten.}

        \subsubsection*{Gas'Danir}
        Gas'Danir gilt bei den anderen Reichen Tirakans als lose Gemeinschaft aus Piraten, Freibeutern und anderem Gesindel. Die inmitten des Felsenmeeres gelegene Insel beherbergt drei größere Städte und viele kleine Siedlungen. Obwohl Gas'Danir keinen guten Ruf hat gibt es sehr wohl eine gesellschaftliche Struktur auf der Insel.

        \abstammungspielwertebox
        {Erhalte +10 auf Heimlichkeit, Infiltration und maritime Bewegung auf kurze Distanz.}
        {In langen offenen Schlagabtauschen erhältst du -10 auf den zweiten und jeden weiteren Nahkampfaustausch derselben Szene.}

        \subsubsection*{Gasdaria}
        Gasdaria ist ein Reich mit einem starken Kastensystem. Das Land, welches reich an Ressourcen wie Gold und exotischen Früchten ist, wird von einer Herrscherkaste geführt. Gasdaria ist das erste Reich welches von den Minotauren auf ihrem Zug nach Norden heimgesucht wurde. Bereits im 3. Jahrhundert gilt das Land größtenteils als Zerstört. Die Bevölkerung ist unter den Hufen versklavt, alle Kontakte zu nördlichen Reichen sind abgebrochen.

        \abstammungspielwertebox
        {Erhalte +10 auf Verwaltung, Logistik und disziplinierte Koordination auf Distanz.}
        {Wenn Pläne unerwartet zusammenbrechen, erhältst du -10 auf die erste Improvisationsprobe.}

        \subsubsection*{Meridian}
        Das Meridianische Reich ist eines der ältesten menschlichen Reiche Tirakans. Es erstreckt sich über die nördliche Landbrücke des Kontinents und ist durch seine Lage an den Küsten des Felsenmeeres und des Ozeans der Stürme ein wichtiger Handelspartner für die Reiche Tirakans. Die Hauptstadt des Reiches ist Merid, eine der ältesten Städte Tirakans und ein bedeutendes Handelszentrum. Meridian hat als eines von zwei Reichen Tirakans eine gewählte Volksvertretung, die alle drei Jahre neu gewählt wird.

        \abstammungspielwertebox
        {Erhalte +10 auf Proben zu Politik, Recht und historischen Präzedenzfällen.}
        {Unter plötzlichem körperlichem Druck (Überfall, Verfolgung) erhältst du -10 auf die erste schnelle Handlungsprobe.}

        \subsubsection*{Nur'Tuk}
        Nur'Tuk ist der Name eines Reitervolkes aus dem Norden Tirakans. Östlich der Grenze Meridians gelegen sind die Steppen der Nur'Tuk, die sich bis zum Gebiet der Katora und der Silkanda erstrecken. Im Laufe der Jahrhunderte kam es immer wieder zu Begegnungen zwischen den Reitern der Nur'Tuk und den Dörfern Meridians. Wenngleich diese zumeist friedlich waren kam es nie zu einem wirklichen Austausch zwischen den beiden Kulturen. Die Nur'Tuk sind kleinwüchsiger als Meridianer, und ihre Haut ist dunkler gefärbt. Sie zählen zu den besten Reitern Tirakans.

        \abstammungspielwertebox
        {Erhalte +10 auf Standhaftigkeitsproben gegen Einschüchterung und Panikdruck.}
        {In Situationen, die Geduld und Feinheit verlangen, erhältst du -10 auf die erste Täuschungs- oder Intrigenprobe.}

        \subsubsection*{Hadewald}
        Hadewald ist das Land der Gilden, der Seefahrt und des Handwerks. In den dichten Wäldern des Reiches ist die Holzwirtschaft das vorherrschende Gewerbe. Anders als in vielen anderen Königreichen herrscht hier nicht das Königshaus, sondern ein Rat der Meister. Der König von Hadewald hat eher repräsentative Aufgaben. Dieser Rat, der von den Handwerksgilden gestellt wird, sorgt für die Gesetzgebung und Rechtsprechung im Land. Die Hadewalder sind für ihren Mut und ihre Stärke bekannt.

        \abstammungspielwertebox
        {Erhalte +10 auf Reitreisen, Tierführung und improvisierte Feldreparaturen.}
        {In höfischen Etiketteszenen erhältst du -10 auf strenge Protokollproben, sofern dir kein anderer SC hilft.}

        \subsubsection*{Toran}
        ``Denke ich an meine Zeit in Toran, so fallen mir zunächst die toraner Tugenden ein. Aufrichtigkeit und Ehrlichkeit sind dem Toraner heilig, ebenso eignet er sich einen außerordentlichen Fleiß und die bekannte Toraner Bescheidenheit an. Das Leben in den Städten ist von regem Treiben geprägt. Überall sieht man bewaffnete, im Sold stehende Männer. Wo die Jugend in Yadosien auf den Feldern arbeitet oder sich dem Müßiggang ergibt, schicken die Toraner ihre Kinder bereits in jungen Jahren auf die Schulbank.'' ~ Jaques Haviere, yadosischer Dichter und Forscher.

        \abstammungspielwertebox
        {Einmal pro Szene erhält eine soziale Probe +10, wenn sie Schulden, Rang oder Politik betrifft.}
        {Verweigerst du eine direkte Herausforderung durch einen Rivalen, kostet dich das 1 Omen (oder +1 Bürde bei 0 Omen).}

        \subsubsection*{Quitaron}
        Die Quitaron sind ein Naturvolk im äußersten Norden Tirakans. Die versträuten Stämme der Quitaron verteilen sich in der weitläufigen Tundra des Nordens. Die Stämme werden von zwei besonders auserkorenen Führern angeführt, einem Wächter und einem Krieger. Das Volk hat sehr wenig Kontakt zu anderen Reichen Tirakans, nur zu den Atiarel Elfen pflegt man feindliche Beziehungen.

        \abstammungspielwertebox
        {Erhalte +10 auf Überleben, Reiten und Ausdauerproben in der Wildnis.}
        {In dichten Stadtvierteln erhältst du -10 auf Orientierung, sofern dich kein ortskundiger Führer begleitet.}

        \subsubsection*{Yadosien}
        Yadosien gilt lange Zeit als das Land der Künste, des Weines und des Müßiggangs. Als direkter Nachfolger des untergegangenen Priesterkaiserreiches Elentrea hat Yadosien eine große Fülle an kulturellem Erbe zu bieten. Zudem ist es über die Jahrhunderte hinweg ein sehr offenes Reich gewesen, in dem jede Kultur, jeder Glaube und auch die Magie gern gesehen waren. All das sollte sich mit einem besonderen Ereignis ändern. Eines Tages wurde ein gefangener Echsenpriester in den Kerker zu Bayard verbracht. Von diesem Tage an sollte in Yadosien nichts mehr so sein, wie es einmal war.

        \abstammungspielwertebox
        {Erhalte +10 auf Proben zu Kunst, Etikette, Diplomatie und dem Erkennen kultureller oder religiöser Zusammenhänge.}
        {Gegen fanatischen Glauben, verbotene Kulte oder verstörende Offenbarungen erhältst du -10 auf die erste Probe zum Widerstehen von Angst, Zweifel oder geistiger Verunsicherung.}

        \subsubsection*{Yavon}
        Die Republik Yavon beherbergt eines der modernsten Staatssysteme Tirakans, wenngleich das Land niemals groß in Erscheinung getreten ist. Still und unauffällig hat das Yavoner Volk eine Revolution losgetreten, ohne von seinen Nachbarn Toran und Meridian beachtet worden zu sein. Auch zu Zeiten des Königshauses hat es niemals Kriege mit den Nachbarn gegeben, lediglich die Orks aus dem Süden sorgten für einige Unruhe. Seit dem 3. Jahrhundert nun ist Yavon eine waschechte Republik. Das Volk wählt einen Rat aus 20 Männern und Frauen, aus jeder Region einen Vertreter.

        \abstammungspielwertebox
        {Erhalte +10 auf Befehls-, Verhandlungs- und gesellschaftliche Proben im Zusammenhang mit Herrschaft.}
        {Wird deine Autorität öffentlich untergraben, erhältst du in diesem Konflikt -10 auf die erste Reaktionsprobe.}

        \subsection{Elfen}\label{subsec:abstammungen-elfen}

        Geschaffen von den Drachen im Zeitalter der Wächter, sind die Elfen das wohl älteste menschenähnliche Volk auf Tirakan. Ihre Hochzeit ist längst vorbei, und heute kann niemand mehr ernsthaft von sich behaupten, jemals einem Elfen begegnet zu sein. Legenden sowie Gelehrte erzählen zwar von einem erneuten Aufleben der Magie, und damit auch der alten Diener, bestätigt ist dies jedoch in keinster Weise.

        \subsubsection*{Ancatir}
        Ancatir steht sowohl für die Elfen des Meeres als auch für den ältesten unsterblichen Sohn der Drachen, der das Volk der Ancatir durch die Jahrhunderte führt. Die Elfen von Ancatir leben abgesondert im Süden des Kontinents und haben sich auf die Seefahrt und das Leben in den kargen Gebieten vor dem Ozean spezialisiert.

        \abstammungspielwertebox
        {Erhalte +10 auf Heimlichkeit und Proben zur Vorbereitung von Fernangriffen aus der Deckung.}
        {Wirst du ohne Vorbereitung in direkten Nahkampf gezwungen, erhält dein erster Nahkampfangriff in jeder Szene -10.}

        \subsubsection*{Atiarel}
        ``\ldots nordwestlich von Meridian hinter dem Tal des Vergessens liegt ein Gebiet, das nur aus Hügeln und Bergen besteht und bis auf eine kleine Stelle im Norden von Wasser umgeben ist. Die östliche Grenze, zu Meridian und den Stammesgebieten der Quitaron, bildet der Strom Kanuba, die restlichen Grenzen das Meer. nur im Norden gibt es eine Landverbindung zu dem Gebiet der Quitaron. Dieses Gebiet besteht aus einem großem Gebirge, daß von Hügeln umschlossen ist. Dieses Gebirge wird Jazabal'Jarh genannt, was in der Sprache der dort lebenden Elfen, den Atiarel, `Heiliges Land' oder `Heiliger Berg' bedeutet. Diese Atiarel sind von großer, schlanker Statur, und von unterschiedlicher Hautfarbe, sei es Normal, wie bei den Menschen, Alabasterweiß oder Steingrau. Sie tragen ihr Haar, welches dünn und seidig ist, entweder lang, als Glatze oder eine, oder beide, Seiten abrasiert. Desweiteren haben sie große, schrägstehende Augen von leuchtender Farbe. Die Atiarel bekommt man recht selten zu Gesicht, da sie ein zurückgelegenes Leben im Gebirge führen \ldots'' ~

        \abstammungspielwertebox
        {Erhalte +10 auf arkane Analyse (Geist + Magiekunde) und das Lesen von Ritualen.}
        {Bei einer katastrophal misslungenen Magieprobe markierst du zusätzlich zum üblichen Rückschlag +1 Bürde.}

        \subsubsection*{Silkanda}
        Das Volk der Silkanda ist eines der 4 Völker der Elfen. Sie sind den Menschen am nächsten, unterhalten seit dem 2. Jahrhundert Beziehungen zu Meridian und Asgoran. Im Laufe der Jahrhunderte nehmen sie mehr und mehr Menschliche Gepflogenheiten an, und akzeptieren bereits im 3. Jahrhundert ihre Währung.

        \abstammungspielwertebox
        {Erhalte +10 auf gesellschaftliche Eleganz (Erscheinung + Einfluss) und diplomatische Erstkontakte.}
        {Wirst du persönlich beleidigt oder entehrt, erhältst du -10 auf Proben zum Lösen oder Rückzug, bis du reagierst.}

        \subsubsection*{Anscharon}
        Anscharon selbst erwachte aus dem magischen Einfluss, und ihm wurde jede Einzelheit des Rituals gewahr. Er zog sich selbst aus Scham und Verzweiflug zurück auf eine Insel jenseits des großen Meeres. Während er apathisch in seiner Verzweiflung lebte, erkannte er nicht, dass Thzularn bereits einen großen Einfluss auf sein Volk genommen hatte. Es entstanden dunkle Kulte, und die Anbetung der Thzularn überwog im Volk der Anscharon. Einzig einige wenige standen zu ihrem Herrn Anscharon.

        \subsection{Zwerge}\label{subsec:abstammungen-zwerge}

        Die Zwergenvölker Tirakans sind höchst unterschiedlich. Die Fraxut Zwerge sind den Menschen wohl gesonnen und eher zugewandt. Dies mag auch daran liegen, dass ihr Gebiet zwischen den Menschlichen Reichen der Yadosier, Hadewalder und Asgoraner liegt. Eher abgeschlagen liegt das Volk der Xordai, welche einen eigenen, sehr starken Glauben haben und eher menschenscheu sind. Die Morgalas hingegen sind eine besondere Art von Zwerg. Sie wurden in grauer Vorzeit von Wisgu versklavt, und folgen sinistren Gelüsten, seit sie vor vielen Jahrhunderten in die Tiefe der Erde gezogen wurden.

        \subsubsection*{Fraxut}
        Die Fraxut sind den Menschen eher zugewandt als andere Zwerge. Dies mag auch der Lage ihres Gebirges geschuldet sein, befindet es sich doch nahe den Reichen der Hadewalder und Yadosier. Die Fraxut ehren ihre vergangenen Könige genauso wie ihre Ahnen, und haben eine sehr starke Gesellschafts- und Glaubensstruktur.

        \abstammungspielwertebox
        {Bei schwachem Licht oder in Dunkelheit erhältst du +10 auf Wahrnehmungs- und Verfolgungsproben.}
        {In hellem offenem Tageslicht erhältst du -10 auf Heimlichkeit, sofern du nicht zusätzliche Zeit für Deckungsvorbereitung aufwendest.}

        \subsubsection*{Xordai}
        Die Xordai sind ein Zwegenvolk im tiefen Süden des Kontinents. Ihrer Geschichtsschreibung nach gibt es sie seit etwa 2500 Jahren, wenngleich die benachbarten O'Grut anderer Ansicht sind. Den Menschen ist nicht viel bekannt über dieses zurückgezogene, ferne Volk, sie sollen erst zum Ende des 9. Jahrhunderts ihre Bekanntschaft machen.

        \abstammungspielwertebox
        {Erhalte +10 auf Bergbau, Steinmetzarbeit und Orientierung unter Tage.}
        {In tiefer Wildnis ohne Straßen oder Steinmarken erhältst du -10 auf Fernorientierung.}

        \subsubsection*{Morgalas}
        Morgalas (frei übersetzt ``Das tiefe Volk'') bezeichnet ein Volk von Zwergen die, obgleich sie zu einem der Ältesten Völker Tirakans gehören, ein tragisches und düsteres Dasein fristen. Der Ursprung der Morgalas liegt fernab jeder bekannten Zivilisation im Osten, jedoch besiedelten sie vor vielen Jahrhunderten einmal den ganzen Kontinent. Im Zeitalter der Diener geschah es, dass der König der Morgalas einen verhängnisvollen Pakt mit Wisgu schloss, der sein Volk auf ewig verfluchen sollte. Mit dem Niedergang der Dämonen zogen sich auch die Morgalas zurück, und waren erst wieder gesehen, als sie auf der Seite von Taurus gegen die Menschen in den Krieg zogen. Die Morgalas werden im Volksmund als Dunkelzwerge bezeichnet.

        \subsection{Orks}\label{subsec:abstammungen-orks}

        Es gibt zwei große Völker der Orks auf Tirakan, die sich jeweils aus einzelnen Stämmen und Stammesverbindungen zusammen setzen. Diese Völker sind in keiner Weise wirklich mit menschlichen Reichen zu vergleichen, es herrschen fortwährend Auseinandersetzungen und kleinere Scharmützel.

        \subsubsection*{Katora}
        Die Sippen der Katora nennen die Länder im südlichen Toran ihr eigen. Die Katora sind nicht wirklich eine feste Gruppe, vielmehr ein Zusammenschluss von einzelnen Sippen. Lediglich zu besonderen Anlässen hört man die Trommeln der Katora schlagen, um grosse Orkzusammenkünfte oder Kriegszüge einzuleiten. Die Katora verheren, wie die Kroto'Chim, den Gott Akhrosch. Akhrosch war es einst, der der erste der Orks war. Noch heute gibt er den Schamenen der Völker ihre Kräfte. Unter den Katora gibt es auch verschieden gesinnte Sippen, die wohl drei grössten sind hier aufgeführt (Insgesamt geht man von gut 20 verschiedenen Hauptsippen aus).

        \abstammungspielwertebox
        {In Dunkelheit erhältst du +10 auf Wahrnehmungsproben und Nahkampfproben.}
        {In hartem direktem Licht erhältst du in jeder Szene -10 auf die erste Präzisionsprobe.}

        \subsubsection*{Kroto'Chim}
        Die Kroto'Chim sind ein Volk der Orks. Sie leben in einem rauhen Land mit den Namen Nuk'Kchtuh, was in ihrer Sprache, dem Chrotok, das fruchtbare, oder das schöne Land heißt. Die Kroto'Chim werden sehr groß und sind sehr robust. In dem rauhen und hügeligem Land, in dem sie leben, haben sich die Kroto'Chim darauf spezialisiert Nahrz'gu und Ziegen zu züchten und Feldarbeit zu betreiben, denn Rinder und edle Feldfrüchte gedeihen hier nicht sehr gut. Der Stolz eines jeden Clans sind ohne Zweifel dessen Burgen, die Menschen duerften sie eher Befestigungsanlagen nennen, und die großen Herden. Die Kroto'Chim verehren, wie die Katora, den Gott Akhrosch, aber auch die den Menschen bekannten Titanen Zyral und Rogal. Akhrosch war es einst, der der erste der Orks war. Noch heute gibt er den Schamenen der Völker ihre Kräfte.

        \abstammungspielwertebox
        {Erhalte +10 auf Menschenkenntnis (Instinkt + Einfluss) und Proben zum Umgang mit Tieren.}
        {In stark urbanen oder künstlichen Umgebungen erhältst du -10 auf Orientierung nach natürlichen Zeichen.}

        \subsection{Goblins}\label{subsec:abstammungen-goblins}

        \subsubsection*{O'Grut}
        Sie sind Leute von kleiner Statur, die im Erwachsenenalter die Größe eines Menschenkindes erreichen. Zugleich kennzeichnet sie eine faltige Haut, die der Haut eines Menschengreises gleicht und von grüngelber Farbe ist. Mit zunehmenden Alter wird die Haut eines O´grut immer dunkler, fast schwarz ist sie bei einem, der das Alter von 50 Jahren erreicht, welches allgemein der Lebenserwartung dieser Wesen entspricht.

        \abstammungspielwertebox
        {Erhalte +10 auf Täuschungsproben und soziale Irreführung, wenn Identität oder Absicht verschleiert sind.}
        {Wirst du öffentlich direkt herausgefordert, erhältst du -10 auf die erste unmittelbare soziale Abwehrprobe, bis du klar Stellung beziehst.}

        \subsubsection*{Asch-Ta-Khi}
        Dieses Goblinvolk, welches auch als die Hochgoblins bezeichnet wird, war ein ursprünglich ein weit verbreitetes Nomadenvolk der Barbarensteppen im 1. und 2. Jahrhundert. Mit dem Erwachen der Stiere zogen sie sich immer weiter zurück in den Norden des Kontinents und verloren ihren stämmischen Zusammenhalt. Es gibt die unterschiedlichsten Stämme und Sippen, überall verteilt auf Tirakan. Neben äusserst kriegerischen, z.T. sogar menschenfressenden Stämmen, gibt es ebenso friedliche oder zurückgezogen lebende.

        \abstammungspielwertebox
        {Erhalte +10 auf Täuschungsproben und soziale Irreführung, wenn Identität oder Absicht verschleiert sind.}
        {Wirst du öffentlich direkt herausgefordert, erhältst du -10 auf die erste unmittelbare soziale Abwehrprobe, bis du klar Stellung beziehst.}

        \subsubsection*{Dunkelgoblins}
        Tief in schattigen Höhlen, Spalten oder auch an den Hängen von Vulkanen ist dieses primitive, wiederwärtige Volk zu finden. Sie messen knapp über einen Schritt, ihre Haut ist faltig und aschgrau. Ihr Maul und ihre Klauen sind übersäht mit spitzen Zähnen und Krallen. Sie tragen meist Felle und führen primitive Keile oder Äxte als Waffen. Manchmal werden sie auch mit verrosteten Kurzschwertern gesehen, die sie wohl von einem ihrer Opfer abgenommen haben. Ihre Sprache besteht nur aus wenigen Worten und gleicht einem krächzendem Grunzen und Schreien, ihre Blutlust ist überall im Land bekannt und vor allem in kleinen abgelegenen Dörfern und Siedlungen gefürchtet. Sie sind feige Geschöpfe der Dunkelheit und wenn sie angreifen, dann nur wenn sie auch zahlenmässig weit überlegen sind. Einige von ihnen hat man schon inmitten des echsischen Heeres, oder als Sklaven von Morgalas oder Asch-Ta-Khi Häuptlingen gesichtet. Ein einzelner Dunkelgoblin ist für einen Krieger kein ernstzunehmender Gegner, doch wo einer von ihnen ist, finden sich meist auch mehrere \ldots

        \subsection{Gnomartige}\label{subsec:abstammungen-gnomartige}

        \subsubsection*{Duigosz}
        Das Land, in dem die Duigoshz leben, ist ein Talkeil, der sich zwischen zwei Gebirgen von Osten nach Westen ins Meer hinein erstreckt. Eine hohe Niederschlagsmenge und der Zufluss aus den Bergen sorgt hier für einen permanent schlammigen Untergrund, sowie für abertausende Wasserrinnsale, die sich in wenigen Metern Abständen nebeneinander anreihend den Weg ins Tal suchen. Die Spitze des Talkeils bildet ein breiter, brüllend reißender Fluss, der noch seinesgleichen in Tirakan sucht.

        \abstammungspielwertebox
        {Erhalte +10 auf Proben für Handwerk, Reparatur oder den Betrieb von Maschinen.}
        {Improvisierst du ohne geeignetes Werkzeug, sinkt der Effekt um eine Stufe (stark zu normal zu begrenzt).}

        \subsubsection*{Gnome}
        Viele Mythen, Legenden und Geschichten ranken sich um dieses wunderliche Völkchen, welches unter der Erde Tirakans haust. Angeblich, so sagen zumindest die Fraxut, fährt die Stadt der Gnome Rolltrutz auf einem riesigen Schienennetz tief unter Erde Tirakans. Vielerorts wird diese Geschichte jedoch als trunkene Zwergengeschichten und Weibergewäsch abgetan.

        \abstammungspielwertebox
        {Erhalte +10 auf Proben für Handwerk, Reparatur oder den Betrieb von Maschinen.}
        {Improvisierst du ohne geeignetes Werkzeug, sinkt der Effekt um eine Stufe (stark zu normal zu begrenzt).}

        \subsubsection*{Gorben}
        Dies ist ein Volk von kleinen zwergenähnlichen Leuten, das allen anderen Völkern gänzlich unbekannt ist. Grund dafür ist sein verstecktes Leben unter Sträuchern und im Unterholz der Wälder. Die Größe von gorbischen Gemeinschaften geht nie über die 12er Grenze und man mag sich wundern, wie die vielen kleinen, über ganz Tirakan verstreuten Gorbensippen den Kontakt zueinander aufrechterhalten. Der Unterschied zwischen Gorben und Zwergen ist, dass erstere noch weitaus kleiner und behaarter sind.

        \abstammungspielwertebox
        {Erhalte +10 auf Proben für Handwerk, Reparatur oder den Betrieb von Maschinen.}
        {Improvisierst du ohne geeignetes Werkzeug, sinkt der Effekt um eine Stufe (stark zu normal zu begrenzt).}

        \subsection{Weitere Kulturen}\label{subsec:abstammungen-weitere}

        \subsubsection*{Doldagor}
        Die Doldagor (ca. 4000 v. EC - 950 EC) sind steinerne Wesen, welche seit dem dritten Zeitalter Tirakan bevölkern. Ihre Gestalt ist die eines Wasserspeiers der Menschen, welche nach dem Bilde der Doldagor gestaltet wurden. Aussehen und Steinart der Doldagor sind höchst unterschiedlich, so können diese Wesen sehr vielgestaltig sein.

        \abstammungspielwertebox
        {Einmal pro Szene gilt vertikale Bewegung für dich als trivial, und du erhältst +10 auf Proben zur Positionierung in der Höhe.}
        {Nach einem misslungenen Manöver mit hoher Beweglichkeit erleidest du eine zusätzliche geringe Konsequenz (Aufprall, Bloßstellung oder Desorientierung).}

        \subsubsection*{Flügler}
        Die Flügler sind kleine, geflügelte Wesen, die große Ähnlichkeit mit den Elfen haben. Ihre Körperproportionen und auch die spitzen Ohren gleichen denen der Elfen. Die Flügler erreichen eine durchschnittliche Größe von ca. 60 bis 70 Fingern. Die Flügel, die ihren Schulterblättern entspringen, können jedoch von Flügler zu Flügler verschieden sein.

        \abstammungspielwertebox
        {Erhalte +10 auf Proben zum Ausweichen oder Umpositionieren in offenem Raum und ignoriere einmal pro Szene ein geringes Bewegungshemmnis.}
        {Bist du gefesselt oder eingeschlossen, beginnst du in dieser Szene mit -10 auf die erste Fluchtprobe.}

        \subsubsection*{Echsen}
        Aus: ``Die Wahrheit über die Welt'' Schriftensammlung, ca. 2000-3000 Jahre alt, unbekannte Autoren (mind. 3) ``Aus der Tiefe des unteren Seins heraus erschuf Aspersia die von ihrem Bilde, welche fortan getrieben würden durch die dunklen Pfade der Unterwelt. Wesen von gleicher Gestalt, die gleichsam ihren Ursprung, wie auch ein Stück von Aspersias Hochmut mit sich tragen. Im Zeichen des Speers mit dem Zeichen der Hörner werden sie kämpfen und es wird Blut über das Land vergossen werden.''

        \subsubsection*{Minotauren}
        Im Jahre 0 erstanden durch Taurus Hand, wüten die Minotauren in den kommenden Jahrhunderten in den südlichen Reichen. Sie nehmen Land um Land, und ihr Feldzug gegen die Menschen überrollt Gasdaria und nähert sich immer weiter dem Felsenmeer. Sie stellen neben den Echsen die größte Gefahr für die Königreiche der Menschen dar.

    \end{multicols*}
}


\section{Wege}\label{sec:wege}

{
    \begin{multicols*}{2}

        \subsection{Weltliche Wege}\label{subsec:wege-weltlich}

        \subsubsection*{Krieger}
        \wegespielwertebox
        {Erhalte +10 auf den ersten Angriff mit einer Nahkampfwaffe in jeder Kampfrunde.}
        {Deine Überzeugung von deinen Fähigkeiten im Kampf lässt dich -10 auf jeden Angriff und jede Verteidung erhalten, wo ein Gefährte neben dir steht.}

        \subsubsection*{Gelehrter}
        \wegespielwertebox
        {Erhalte +10 auf Geschichte, Mythen und Proben zur Zusammenführung von Erkenntnissen.}
        {Bei plötzlicher Gewalt erhältst du in dem Konflikt -10 auf die erste Probe.}

        \subsubsection*{Abenteurer}
        \wegespielwertebox
        {Einmal pro Szene ignorierst du ein geringes Hindernis der Fortbewegung und erhältst +10 auf Proben zu Erkundungsrisiken.}
        {Bist du überdehnt (wenig Vorräte oder keine Rast), markierst du nach einer misslungenen Gefahrenprobe +1 Bürde.}

        \subsubsection*{Barde}
        \wegespielwertebox
        {Erhalte +10 auf Auftritte, soziale Einflussnahme und Szenen, in denen du Moral stärkst.}
        {Wirst du öffentlich verspottet oder diskreditiert, erhältst du danach -10 auf die erste Überzeugungsprobe.}

        \subsubsection*{Dieb}
        \wegespielwertebox
        {Erhalte +10 auf heimliches Eindringen, Schlossöffnung und Proben auf Fingerfertigkeit.}
        {Nach einer misslungenen Heimlichkeitsprobe erhältst du in derselben Szene -10 auf deine nächste soziale Tarnprobe.}

        \subsubsection*{Händler}
        \wegespielwertebox
        {Erhalte +10 auf Feilschen, Wertschätzung und Proben auf Handelsnetzwerke.}
        {In Geschenkpflicht- oder Tabuökonomien erhältst du -10 auf die erste Handelsverhandlungsprobe.}

        \subsubsection*{Mechanikus}
        \wegespielwertebox
        {Erhalte +10 auf Proben die Schätzen, Rechnen oder Baumeisterei enthalten.}
        {Durch ständiges Grübeln oder Bewerten der Welt in einem technischen Sinn erhälst du -10 auf alle Proben, die Kommunikation enthalten.}

        \subsubsection*{Seefahrer}
        \wegespielwertebox
        {Erhalte +10 auf Seefahrt, Schiffskunde und Navigationsproben auf See.}
        {Bei Überlandreisen durch Berge oder Wüste erhältst du an jedem Tag -10 auf die erste Orientierungsprobe.}

        \subsubsection*{Pirat}
        \wegespielwertebox
        {Erhalte +10 auf Seefahrt, Kommunikationsproben und Navigationsproben auf See.}
        {Bei Überlandreisen erhältst du an jedem Tag -10 auf die erste Orientierungsprobe.}

        \subsubsection*{Söldner}
        \wegespielwertebox
        {Erhalte +10 bei Vertragsverhandlungen, Drohgebärden und praktischer Scharmützeltaktik.}
        {Werden Bezahlung, Bedingungen oder Loyalität öffentlich infrage gestellt, erhältst du -10 auf die erste Probe.}

        \subsubsection*{Wegelagerer}
        \wegespielwertebox
        {Erhalte +10 beim Fallen stellen, wenn du aus dem Hinterhalt agierst oder dich im Verborgenen bewegst.}
        {Bei allen sozialen Proben in Gesellschaft erhälst du -10, bei gehobener Gesellschaft -20.}

        \subsubsection*{Meuchler}
        \wegespielwertebox
        {Erhalte +10 auf geplante Tötungsversuche und Proben zur Vorbereitung durch Zielstudium.}
        {In langem offenem Kampf erhältst du ab der zweiten Runde -10 auf Angriffe.}

        \subsubsection*{Gaukler}
        \wegespielwertebox
        {Erhalte +10 auf Proben die Darstellung, Täuschen oder auch Betören.}
        {Du bist anfällig für Aberglauben, erhalte -10 auf alle Proben, die rationales Denken erfordern.}

        \subsubsection*{Medikus}
        \wegespielwertebox
        {Erhalte +10 auf Proben auf Heilung, Erste Hilfe und Naturkunde.}
        {Erhalte -10 auf alle Proben die seichte Kommunikation erfordern.}

        \subsubsection*{Jäger}
        \wegespielwertebox
        {Erhalte +10 auf Verfolgung, die Vorbereitung von Schüssen und das Lesen von Spuren in Wildnisszenen.}
        {In höfischen oder juristischen Verhandlungsszenen erhältst du -10 auf die erste Einflussprobe.}

        \subsubsection*{Spion}
        \wegespielwertebox
        {Erhalte +10 auf Infiltration, Überwachung und verdeckte Kommunikationsproben.}
        {Ein bekannter Rivale kann dir -10 auf die erste soziale Probe einer Szene geben, in der deine Identität offengelegt ist.}

        \subsubsection*{Soldat}
        \wegespielwertebox
        {Erhalte +10 auf koordinierte Kampfmanöver und Proben auf Schlachtfelddisziplin.}
        {Bei direkten widersprüchlichen Befehlen oder Befehlszusammenbruch markierst du +1 Bürde oder erhältst in dieser Runde -10 auf Proben, die über deine Handlungsreihenfolge entscheiden.}

        \subsubsection*{Geistlicher}
        \wegespielwertebox
        {Erhalte +10 auf Riten, Liturgie und Überzeugungsproben aus der Lehre heraus.}
        {Handelst du offen gegen deine erklärte Lehre, verlierst du 1 Omen (oder markierst +1 Bürde bei 0 Omen).}

        \subsubsection*{Schmied}
        \wegespielwertebox
        {Erhalte +10 auf Schmieden, Reparatur und Proben zur Beurteilung von Metall.}
        {Ohne geeignetes Werkzeug oder Schmiede sinkt der Effekt bei Ausrüstungsreparaturen um eine Stufe.}

        \subsubsection*{Boxer}
        \wegespielwertebox
        {Erhalte +10 auf unbewaffneten Nahkampf und Proben auf Schmerzausdauer.}
        {Gegen Reichweitenwaffen auf offenem Boden erhältst du in jeder Szene -10 auf den ersten Nahkampfaustausch.}

        \subsubsection*{Schreiber}
        \wegespielwertebox
        {Erhalte +10 auf Aufzeichnungen, Schriften und Proben zur Echtheit von Dokumenten.}
        {Unter unmittelbarem Zeitdruck erhältst du in der Szene -10 auf die erste Schreib- oder Analyseprobe.}

        \subsubsection*{Bote}
        \wegespielwertebox
        {Erhalte +10 auf schnelle Reisen und Proben auf Eilzustellung unter Druck.}
        {Bist du mit schwerer Ausrüstung beladen, erhältst du -10 auf Verfolgungs- oder Fluchtbewegung.}

        \subsubsection*{Bestatter}
        \wegespielwertebox
        {Erhalte +10 auf Totenriten, Leichenumgang und ruhiges Handeln in Szenen des Todes.}
        {Lebende Menschen misstrauen deinem Ton; in festlichen Situationen erhältst du -10 auf die erste charmante soziale Probe.}

        \subsubsection*{Ritter}
        \wegespielwertebox
        {Erhalte +10 auf berittene Kampfkommandos, Ehrenduelle und formale Autoritätshaltung.}
        {Das Ignorieren einer erklärten Herausforderung kostet 1 Omen (oder +1 Bürde bei 0 Omen).}

        \subsubsection*{Medium}
        \wegespielwertebox
        {Erhalte +10, wenn du Omen, Echos oder Anzeichen geisterhafter Präsenz deutest.}
        {Bei katastrophal misslungenem okkultem Kontakt markierst du zusätzlich zum normalen Nachhall +1 Bürde.}

        \subsubsection*{Bader}
        \wegespielwertebox
        {Erhalte +10 auf grobe Traumaversorgung und Proben zur Stabilisierung im Feld.}
        {Bei feiner langwieriger Behandlung erhältst du -10 auf die erste präzise medizinische Probe.}

        \subsubsection*{Abdecker}
        \wegespielwertebox
        {Erhalte +10 auf Einschüchterung, die Verarbeitung von Kadavern und Ausdauer bei grimmiger Arbeit.}
        {In etikettelastigen sozialen Szenen erhältst du -10 auf die erste Diplomatieprobe.}

        \subsubsection*{Wirt}
        \wegespielwertebox
        {Erhalte +10 auf Gerüchtejagd, das Lesen von Menschenmengen und Proben zur Gastfreundschaftslogistik.}
        {Beginnt eine Szene mit offener Gewalt, erhältst du -10 auf die erste Befehls- oder Deeskalationsprobe.}

        \subsection{Magische Wege}\label{subsec:wege-magisch}

        \subsubsection*{Chimärologe}
        \wegespielwertebox
        {Erhalte +10 auf Proben, die Anatomie, Alchemie oder das Verändern lebender Materie betreffen.}
        {Misslingt dir eine Probe auf lebende Experimente oder körperliche Veränderung katastrophal, markierst du zusätzlich +1 Bürde.}

        \subsubsection*{Dämonologe}
        \wegespielwertebox
        {Erhalte +10 auf Proben, die Beschwörung, Bannung oder das Erkennen dämonischer Einflüsse betreffen.}
        {Bei katastrophal misslungener dämonischer Magie oder Kontaktaufnahme markierst du zusätzlich +1 Bürde.}

        \subsubsection*{Druide}
        \wegespielwertebox
        {Erhalte +10 auf Proben, die Naturmagie, Pflanzenkunde oder das Deuten wilder Zeichen betreffen.}
        {In stark urbanen oder verdorbenen Umgebungen erhältst du -10 auf die erste Probe mit Naturmagie oder Naturdeutung in der Szene.}

        \subsubsection*{Geisterbeschwörer}
        \wegespielwertebox
        {Erhalte +10 auf Proben, die Anrufung, Bindung oder Deutung von Geistern und Echos betreffen.}
        {Bei katastrophal misslungener Geisterbeschwörung oder Kontaktaufnahme markierst du zusätzlich +1 Bürde.}

        \subsubsection*{Hexer}
        \wegespielwertebox
        {Erhalte +10 auf Proben, die Flüche, Hexereien oder das Wirken durch verborgene Pakte betreffen.}
        {Wird eine Hexerei von dir katastrophal misslungen oder gegen dich zurückgeworfen, markierst du zusätzlich +1 Bürde.}

        \subsubsection*{Hermetiker}
        \wegespielwertebox
        {Erhalte +10 auf Proben, die Formeln, Siegel, arkanes Ordnen oder kontrollierte Ritualmagie betreffen.}
        {Wirst du bei hermetischer Magie aus deinem vorbereiteten Ablauf gerissen, erhältst du -10 auf die erste Folgeprobe derselben Szene.}

        \subsubsection*{Mystiker}
        \wegespielwertebox
        {Erhalte +10 auf Proben, die Visionen, innere Einkehr oder das Deuten von Omen und verborgenen Zusammenhängen betreffen.}
        {Wirst du aus Trance, Kontemplation oder innerer Sammlung jäh gerissen, erhältst du -10 auf die erste Probe geistiger Klarheit in der Szene.}

        \subsubsection*{Nekrologe}
        \wegespielwertebox
        {Erhalte +10 auf Proben, die Totenriten, nekromantische Erkenntnis oder den Umgang mit sterblichen Überresten betreffen.}
        {In Szenen mit lebendiger Festlichkeit, Fruchtbarkeit oder ausgelassener Gemeinschaft erhältst du -10 auf die erste soziale Einflussprobe.}

        \subsubsection*{Runenleger}
        \wegespielwertebox
        {Erhalte +10 auf Proben, die Runeninschriften, magische Zeichen oder die Vorbereitung gebundener Wirkungen betreffen.}
        {Musst du ohne Zeit für saubere Vorbereitung oder ohne geeignete Fläche wirken, erhältst du -10 auf die erste passende Probe der Szene.}

        \subsubsection*{Schwarzmagier}
        \wegespielwertebox
        {Erhalte +10 auf Proben, die zerstörerische, verderbende oder furchterregende Magie betreffen.}
        {Zeigst du in einer Szene offen schwarze Magie, erhältst du -10 auf die erste anschließende soziale Probe gegen Zeugen oder Betroffene.}

        \subsubsection*{Weissmagier}
        \wegespielwertebox
        {Erhalte +10 auf Proben, die Heilung, Schutz oder reinigende Magie betreffen.}
        {Verursachst du vorsätzlich Leid oder nutzt deine Magie offen zur Grausamkeit, erhältst du -10 auf die nächste passende Magieprobe.}

        \subsubsection*{Zauberer}
        \wegespielwertebox
        {Erhalte +10 auf Proben, die freie arkane Zauberei, magische Improvisation oder das Anpassen bekannter Wirkungen betreffen.}
        {Wirkst du unter starkem unmittelbarem Druck ohne Vorbereitung, erhältst du -10 auf die erste Zauberprobe der Szene.}

        \subsubsection*{Waldläufer}
        \wegespielwertebox
        {Erhalte +10 auf Verfolgung, Feldhandwerk und die Vorbereitung von Fernhaltungen im wilden Gelände.}
        {In dichten Stadtvierteln erhältst du -10 auf Orientierung, sofern dir kein ortskundiger Führer hilft.}

        \subsection{Wege der Götter}\label{subsec:wege-goetter}

        \subsubsection*{Mönch}
        \wegespielwertebox
        {Erhalte +10 auf Proben, die Askese, innere Disziplin oder waffenlose Selbstbeherrschung betreffen.}
        {In Szenen des Überflusses, der Ausschweifung oder starker weltlicher Versuchung erhältst du -10 auf die erste soziale oder geistige Widerstandsprobe.}

        \subsubsection*{Paladin}
        \wegespielwertebox
        {Erhalte +10 auf Proben, die heiligen Kampf, Schutz Unschuldiger oder standhafte Glaubensautorität betreffen.}
        {Weichst du in einer Szene offen von deinem Eid oder deiner heiligen Pflicht ab, erhältst du -10 auf die nächste passende Glaubens- oder Standhaftigkeitsprobe.}

        \subsubsection*{Priester}
        \wegespielwertebox
        {Erhalte +10 auf Proben, die Liturgie, Seelsorge oder das Anrufen göttlicher Ordnung betreffen.}
        {Wirst du in einer Szene öffentlich der Heuchelei, Glaubensschwäche oder Entweihung beschuldigt, erhältst du -10 auf die erste Überzeugungsprobe.}

        \subsubsection*{Vampirjäger}
        \wegespielwertebox
        {Erhalte +10 auf Proben, die das Aufspüren, Erkennen oder Bekämpfen von Vampiren und blutgebundenen Schrecken betreffen.}
        {In Gegenwart offensichtlicher vampirischer Verderbnis erhältst du -10 auf die erste Probe zu Zurückhaltung, Heimlichkeit oder diplomatischem Abwarten.}

    \end{multicols*}
}


\section{Bindungen}\label{sec:bindungen}

\begin{multicols*}{2}

    \noindent\textbf{Verfallene Kapelle}\par
    \begin{itemize}[leftmargin=*,nosep]
\item Vorteil: Erhalte +10 auf Proben mit Ritus, Mythen oder Standhaftigkeit, die mit geweihten Ruinen, Eiden oder Begräbnisriten zusammenhängen.
\item Verwundbarkeit: Wenn du in einem heiligen Ort die Hilflosen zurücklässt, markierst du +1 Bürde und die Bindung ist erschöpft.
    \end{itemize}
    \noindent\textbf{Schwarzsalz-Kompanie}\par
    \begin{itemize}[leftmargin=*,nosep]
\item Vorteil: Erhalte +10 auf Vertragshebel, Söldnerverhandlungen oder Schlachtfeldkoordination mit Veteranen.
\item Verwundbarkeit: Verweigerst du einem früheren Mitglied der Kompanie unter Druck Hilfe, verlierst du 1 Omen (oder markierst +1 Bürde bei 0 Omen).
    \end{itemize}
    \noindent\textbf{Trauerschwester}\par
    \begin{itemize}[leftmargin=*,nosep]
\item Vorteil: Einmal pro Szene kannst du eine soziale Konsequenz oder Druckkonsequenz um eine Stufe senken, wenn du handelst, um Verwandte zu schützen.
\item Verwundbarkeit: Wird deine Schwester bedroht, erhält deine erste Handlung ohne Unterstützungsabsicht in dieser Szene -10.
    \end{itemize}
    \noindent\textbf{Alter Mentor}\par
    \begin{itemize}[leftmargin=*,nosep]
\item Vorteil: Erhalte einmal pro Szene +10 auf eine Vorbereitungs-, Recherche- oder taktische Aufbauprobe, nachdem du die Lehre deines Mentors benannt hast.
\item Verwundbarkeit: Ignorierst du im Spiel eine bekannte Warnung deines Mentors, erhält deine nächste passende Probe -10.
    \end{itemize}
    \noindent\textbf{Stadtunterkrypta}\par
    \begin{itemize}[leftmargin=*,nosep]
\item Vorteil: Erhalte +10 auf Heimlichkeit, Orientierung und okkultes Erinnern in Katakomben, Abwasserläufen und Grabgängen.
\item Verwundbarkeit: In offenen Tageslichtvierteln erhältst du in jeder Szene -10 auf die erste Wahrnehmungs- oder Orientierungsprobe.
    \end{itemize}
\end{multicols*}


\section{Male}\label{sec:male}

\begin{multicols*}{2}

    \noindent\textbf{Verbrannte Hände}\par
    \begin{itemize}[leftmargin=*,nosep]
\item Vorteil: Erhalte +10 auf Proben zu Schmerzausdauer und Hitzegefahren.
\item Verwundbarkeit: Bei feinen Handarbeiten erhältst du -10, sofern du nicht geeignetes Werkzeug oder zusätzliche Zeit einsetzt.
    \end{itemize}
    \noindent\textbf{Eidmal}\par
    \begin{itemize}[leftmargin=*,nosep]
\item Vorteil: Erhalte +10 auf Befehl oder Einschüchterung, wenn du dein geschworenes Mal anrufst.
\item Verwundbarkeit: Wenn du dein gesprochenes Versprechen in der Szene brichst, kostet dich das 1 Omen (oder +1 Bürde bei 0 Omen).
    \end{itemize}
    \noindent\textbf{Fehlendes Auge}\par
    \begin{itemize}[leftmargin=*,nosep]
\item Vorteil: Erhalte +10 auf das Lesen von Hinterhaltswinkeln und das Vorausahnen von Bedrohungen.
\item Verwundbarkeit: Bei Präzisionsproben auf lange Distanz erhältst du -10.
    \end{itemize}
    \noindent\textbf{Gebrochener Kiefer}\par
    \begin{itemize}[leftmargin=*,nosep]
\item Vorteil: Erhalte +10 auf das Widerstehen von Zwang, Panik oder Verhördruck.
\item Verwundbarkeit: In einer formalen Szene erhältst du -10 auf die erste Charme- oder Diplomatieprobe.
    \end{itemize}
    \noindent\textbf{Nachtzittern}\par
    \begin{itemize}[leftmargin=*,nosep]
\item Vorteil: Erhalte +10, wenn du im Lager oder in der Ruhe Gefahr erspürst.
\item Verwundbarkeit: Nach einer Nacht ohne Sicherheit beginnst du die nächste Szene mit +1 Bürde.
    \end{itemize}
    \noindent\textbf{Pockennarbige Haut}\par
    \begin{itemize}[leftmargin=*,nosep]
\item Vorteil: Erhalte +10 auf Einschüchterung gegen Abergläubische.
\item Verwundbarkeit: Bei Etikette oder höfischen Ersteindrücken erhältst du -10.
    \end{itemize}
    \noindent\textbf{Hinkendes Bein}\par
    \begin{itemize}[leftmargin=*,nosep]
\item Vorteil: Erhalte +10 auf Standfestigkeit, wenn du unter Druck Boden hältst.
\item Verwundbarkeit: Bei Verfolgungen oder schnellem Umpositionieren erhältst du -10 auf die erste Bewegungsprobe.
    \end{itemize}
    \noindent\textbf{Zerschlagene Knöchel}\par
    \begin{itemize}[leftmargin=*,nosep]
\item Vorteil: Erhalte +10 auf unbewaffnete Schläge und Schmerzwiderstand.
\item Verwundbarkeit: Bei feiner Handwerks- oder Heilarbeit sinkt der Effekt um eine Stufe.
    \end{itemize}
    \noindent\textbf{Rituelle Narbung}\par
    \begin{itemize}[leftmargin=*,nosep]
\item Vorteil: Erhalte +10 auf okkulte Wiedererkennung und das Erinnern von Riten.
\item Verwundbarkeit: Unheilvolle Zeichen und Gesänge können dich ablenken: Die erste Konzentrationsprobe erhält -10.
    \end{itemize}
    \noindent\textbf{Würgemal am Hals}\par
    \begin{itemize}[leftmargin=*,nosep]
\item Vorteil: Erhalte +10 auf Entfesselungskunst und Proben zum Brechen von Fesseln.
\item Verwundbarkeit: In Szenen mit erstickendem Rauch, Würgegas oder tiefem Wasser erhältst du -10 auf die erste Ausdauerprobe.
    \end{itemize}
    \noindent\textbf{Versengte Lungen}\par
    \begin{itemize}[leftmargin=*,nosep]
\item Vorteil: Erhalte +10 auf das Widerstehen von Rauch- und Ascheeinwirkung.
\item Verwundbarkeit: Sprinten oder lange Anstrengung geben in einer Szene -10 auf die zweite und jede weitere körperliche Probe.
    \end{itemize}
    \noindent\textbf{Verstümmeltes Ohr}\par
    \begin{itemize}[leftmargin=*,nosep]
\item Vorteil: Erhalte +10 auf Lippenlesen und das Deuten von Körpersprache in Menschenmengen.
\item Verwundbarkeit: In lärmendem Chaos mit schlechter Sicht erhältst du -10 auf hörbasierte Wahrnehmung.
    \end{itemize}
    \noindent\textbf{Angebrochene Rippen}\par
    \begin{itemize}[leftmargin=*,nosep]
\item Vorteil: Erhalte +10 darauf, stumpfe Gewalt zu überstehen.
\item Verwundbarkeit: Beim Klettern, Ringen oder schweren Heben erhältst du -10, sofern du dich nicht abstützt.
    \end{itemize}
    \noindent\textbf{Eisenstift in der Schulter}\par
    \begin{itemize}[leftmargin=*,nosep]
\item Vorteil: Erhalte +10 auf Waffenhalt und das Widerstehen gegen Entwaffnung.
\item Verwundbarkeit: Schläge oder Würfe über Kopf erhalten ohne Vorbereitung -10.
    \end{itemize}
    \noindent\textbf{Erfrorene Finger}\par
    \begin{itemize}[leftmargin=*,nosep]
\item Vorteil: Erhalte +10 auf Überlebensproben in extremer Kälte.
\item Verwundbarkeit: In nasskalten Bedingungen erhältst du -10 auf die erste präzise Handprobe.
    \end{itemize}
    \noindent\textbf{Ausgebrannte Wunde}\par
    \begin{itemize}[leftmargin=*,nosep]
\item Vorteil: Einmal pro Szene ignorierst du eine geringe Blutungskonsequenz.
\item Verwundbarkeit: Magische Heilung an dir ist instabil: Eine misslungene Heilprobe fügt eine zusätzliche Komplikation hinzu.
    \end{itemize}
    \noindent\textbf{Vom Peitschenhieb gezeichneter Rücken}\par
    \begin{itemize}[leftmargin=*,nosep]
\item Vorteil: Erhalte +10 auf Gewaltmärsche und das Tragen von Last.
\item Verwundbarkeit: Plötzliche Treffer auf den Oberkörper geben -10 auf unmittelbare Reaktionsproben.
    \end{itemize}
    \noindent\textbf{Salz in den Adern}\par
    \begin{itemize}[leftmargin=*,nosep]
\item Vorteil: Erhalte +10 auf das Widerstehen von Gift und Krankheit.
\item Verwundbarkeit: Starke alchemische oder medizinische Geschmäcker können Würgereiz auslösen: Die erste Probe auf eingenommene Mittel erhält -10.
    \end{itemize}
    \noindent\textbf{Überlebende des Grabfiebers}\par
    \begin{itemize}[leftmargin=*,nosep]
\item Vorteil: Erhalte +10 auf Standhaftigkeitsproben gegen Untote und Leichenbedrohungen.
\item Verwundbarkeit: In der Nähe offener Gräber oder Beinhaufen markierst du +1 Bürde, sofern dir nicht Körper + Standhaftigkeit gelingt.
    \end{itemize}
    \noindent\textbf{Gespaltene Zungennarbe}\par
    \begin{itemize}[leftmargin=*,nosep]
\item Vorteil: Erhalte +10 auf Codesprache, Gaunersprache und rituelle Formulierungen.
\item Verwundbarkeit: In formaler Rhetorik und öffentlicher Rede erhältst du -10 auf die erste Überzeugungsprobe.
    \end{itemize}
\end{multicols*}
